% TOP_PROBLEM: {"problemId": "problem.nirenberg-problem-on-the-two-sphere", "canonicalTitle": "Nirenberg Problem on the Two-Sphere", "releaseRank": 328, "primaryDomain": "geometry_topology", "primaryDomainLabel": "Geometry & topology", "displayStatus": "open_with_solved_subcases", "releaseStatus": "open_with_solved_subcases"}
% !TeX program = lualatex
% ATTEMPT_STATUS: unresolved
% Model: GPT-6 Astra Ultra; continuation dated 27 September 2026.
\documentclass[11pt]{article}
\usepackage[margin=25mm]{geometry}
\usepackage{fontspec}
\setmainfont{Latin Modern Roman}
\usepackage{amsmath,amssymb,mathtools}
\usepackage{unicode-math}
\setmathfont{Latin Modern Math}
\usepackage{xurl,hyperref}
\hypersetup{colorlinks=true,urlcolor=blue}
\setcounter{secnumdepth}{0}
\setlength{\parindent}{0pt}
\setlength{\parskip}{5pt}
\setlength{\emergencystretch}{3em}
\begin{document}
\section*{\texorpdfstring{Nirenberg Problem on the Two-Sphere}{Nirenberg Problem on the Two-Sphere}}\label{328}
\subsection{Definitions and mathematical statement}
Let $g_0$ be the unit round metric on $S^2$, with Gaussian curvature one, and take the Laplacian convention $\Delta_{g_0}=\operatorname{div}\nabla$. A smooth conformal metric is $g=e^{2u}g_0$ with $u\in C^\infty(S^2,{\mathbb{R}})$. Its curvature is
\[
 K_g=e^{-2u}(1-\Delta_{g_0}u).
\]
The Nirenberg problem asks for necessary and sufficient conditions on $K\in C^\infty(S^2,{\mathbb{R}})$ for solvability of
\[
 -\Delta_{g_0}u+1=K e^{2u}\quad\text{on }S^2.
\]
In particular, a solution must satisfy $\int_{S^2}K e^{2u}{\,\mathrm{d}} A_{g_0}=4\pi$, but this identity involves the unknown $u$ and is not a complete characterization. The task is exact existence within the fixed conformal class, not uniqueness of $u$ or existence of a curvature function after an unrelated change of metric.
\par\smallskip\textit{Formulation and concept references:} \href{https://link.springer.com/article/10.1007/s00032-022-00370-1}{[S1]}.

\subsection{Short English statement}
Which smooth functions on the sphere occur as Gaussian curvature after a smooth conformal change of the round metric?
\subsection{Research attempt}
\paragraph{Scope and current references.}
A complete necessary-and-sufficient characterization of smooth
prescribed curvature is \textbf{UNRESOLVED} here. The source [S1]
discusses bubbling and loss of compactness. The newer reference [S2],
submitted 31 March 2026, reports verified existence near particular
computed approximations and determines their symmetries. Its abstract
does not claim a characterization for every smooth $K$. We derive
an exact obstruction, prove a local existence class, construct
nonconstant examples, and show explicitly why a naive compactness
argument fails even for constant positive curvature.

All integrations in this attempt use the round area form $dA_0$.
The sign convention is $\Delta=\operatorname{div}\nabla$, so
$\int v\Delta v=-\int|\nabla v|^2$ and the coordinate functions
$x_i$ on the unit sphere satisfy $\Delta x_i=-2x_i$.

\paragraph{Elementary necessary conditions and constant data.}
Integrating the equation gives
\[
 \int_{S^2}K e^{2u}\,dA_0=4\pi.
 \tag{328.1}
\]
Therefore $K$ must be positive somewhere. This excludes nonpositive
functions, including zero, but it is not sufficient. If $K\equiv c>0$,
the constant $u=-\tfrac12\log c$ is a solution. If $K$ is multiplied
by a positive constant $a$, replacing a solution $u$ by
$u-\tfrac12\log a$ supplies a solution for $aK$. This scaling
changes area and keeps the equation's constant term equal to one.

At a maximum of $u$, the equation gives $K e^{2u}\geq1$ at that
point. At a minimum it gives $K e^{2u}\leq1$. These are statements
at extrema of the unknown $u$, not bounds on the extrema of $K$
alone. A maximum-principle argument that identifies their locations
would introduce an unproved relation.

\paragraph{Deriving the conformal-vector-field obstruction.}
Fix a coordinate function $z=x_i$, and let $X=\nabla z$. On the
round sphere,
\[
 \operatorname{div}X=-2z,\qquad
 \operatorname{Hess}z=-z g_0.
 \tag{328.2}
\]
These identities follow, for example, by differentiating the
restriction of a linear Euclidean coordinate to the sphere. We prove
the Kazdan--Warner identity directly rather than treating it as an
extra hypothesis:
\[
 \int_{S^2}X(K)e^{2u}\,dA_0=0.
 \tag{328.3}
\]
Differentiate $Ke^{2u}=1-\Delta u$ in the $X$ direction. The left
side of (328.3) becomes
\[
 -\int X(\Delta u)-2\int X(u)+2\int X(u)\Delta u.
 \tag{328.4}
\]
Integration by parts using (328.2) gives
$\int X(\Delta u)=2\int z\Delta u=-4\int zu$ and
$\int X(u)=2\int zu$, so the first two terms cancel.
For the last term,
\[
 \begin{split}
 \int X(u)\Delta u
 &=-\int\operatorname{Hess}u(X,\nabla u)
          -\int\operatorname{Hess}z(\nabla u,\nabla u)\\
 &=-\tfrac12\int X(|\nabla u|^2)+\int z|\nabla u|^2=0.
 \end{split}
\]
The final equality again uses $\operatorname{div}X=-2z$. This proves
(328.3), including its weight $e^{2u}$. Removing that weight would
give a different and generally false necessary condition.

\paragraph{Strictly positive curvature data can still be impossible.}
Take $K=a+bz$ with $b\ne0$. Since $|\nabla z|^2=1-z^2$,
\[
 X(K)=b(1-z^2).
\]
It has a strict, fixed sign away from the two poles. Because
$e^{2u}$ is everywhere positive, its integral in (328.3) cannot
vanish. Thus no smooth conformal solution exists. In particular,
$a>|b|>0$ gives an everywhere positive function which is not
realizable in the fixed conformal class.

This is an obstruction for the exact function at its given points.
Changing coordinates by a conformal transformation acts on the
curvature function by composition and also changes the conformal
factor. It cannot be used to replace the desired $K$ with an
unrelated function and still claim to have solved the original
prescription.

\paragraph{A local existence theorem for antipodally even data.}
Let $0<\alpha<1$ and use the Banach spaces of antipodally even
real functions $C^{2,\alpha}_{\mathrm{even}}$ and
$C^{0,\alpha}_{\mathrm{even}}$. Define
\[
 \mathcal F(u)=e^{-2u}(1-\Delta u).
\]
It is a smooth map between these spaces, $\mathcal F(0)=1$, and
\[
 D\mathcal F_0(v)=-\Delta v-2v.
 \tag{328.5}
\]
The spherical harmonic eigenvalues of $-\Delta$ are
$\ell(\ell+1)$, and the eigenspace for $\ell$ has antipodal
parity $(-1)^\ell$. The only zero eigenvalue of (328.5) occurs
at $\ell=1$, which is odd. Thus (328.5) has trivial kernel on
the even subspace and is an isomorphism there. More explicitly,
self-adjoint elliptic Fredholm theory gives surjectivity after the
kernel is removed, and the Schauder estimate gives a bounded inverse.
These standard elliptic and spectral facts are the external inputs.

The Banach-space inverse-function theorem now gives a neighborhood
of $1$ in $C^{0,\alpha}_{\mathrm{even}}$ in which every $K$ has
a unique sufficiently small even $C^{2,\alpha}$ solution. For
smooth $K$, the equation and Schauder bootstrapping make $u$ smooth.
This proves a genuine local prescription result with a clearly
stated symmetry and norm restriction.

On the full function space (328.5) has its three-dimensional
coordinate-function kernel. Ignoring it would make an inverse-function
argument invalid. Removing that kernel by a conformal gauge requires
compatibility with the nonlinear obstruction (328.3), not merely
discarding three equations. The even subspace avoids it because
the symmetry excludes exactly those modes; it does not imply
global solvability for all even positive functions.

\paragraph{A directly manufactured nonconstant family.}
Put $P_2(z)=(3z^2-1)/2$, so $\Delta P_2=-6P_2$. For real
$\varepsilon$ let
\[
 u_\varepsilon=\varepsilon P_2(z),\qquad
 K_\varepsilon=e^{-2\varepsilon P_2(z)}
                           (1+6\varepsilon P_2(z)).
 \tag{328.6}
\]
Substitution proves the equation exactly. Since $-1/2\leq P_2\leq1$,
the curvature is strictly positive whenever
$-1/6<\varepsilon<1/3$. It is nonconstant for nonzero sufficiently
small $\varepsilon$, as follows either by differentiating at zero
or by comparing its values as a function of $P_2$.
Moreover $\int P_2\,dA_0=0$, so
$\int K_\varepsilon e^{2u_\varepsilon}\,dA_0=4\pi$ directly.
The identity (328.3) also holds: for $z=x_3$, its integrand is odd
under $z\mapsto-z$, while for the other coordinate directions
reflection symmetry in those coordinates gives the same cancellation.

Manufacturing $K$ from $u$ verifies a class of examples but does not
invert the curvature map for arbitrary prescribed data. In particular,
it supplies no criterion deciding whether an independently given
$K$ belongs to the image of that map.

\paragraph{An explicit failure of compactness with $K=1$.}
In stereographic coordinates $w\in\mathbb R^2$ the round metric is
$4|dw|^2/(1+|w|^2)^2$. Pull it back under the conformal dilation
$w\mapsto\lambda w$. The resulting conformal factor is
\[
 u_\lambda(w)=\log\lambda+\log(1+|w|^2)
                         -\log(1+\lambda^2|w|^2).
 \tag{328.7}
\]
Each factor extends smoothly over the point at infinity, because it
comes from a smooth conformal automorphism of the sphere. Pullback
preserves constant Gaussian curvature, so every $u_\lambda$ solves
the equation for $K=1$. Yet $u_\lambda(0)=\log\lambda\to\infty$.

The corresponding area measure is
\[
 e^{2u_\lambda}dA_0
       =\frac{4\lambda^2}{(1+\lambda^2|w|^2)^2}\,dw.
\]
Its total mass is $4\pi$, while its mass outside the coordinate
disk of radius $R>0$ is exactly
\[
 \frac{4\pi}{1+\lambda^2R^2}\longrightarrow0.
 \tag{328.8}
\]
This follows by one radial integral. Thus a fixed total curvature
and positive smooth prescribed curvature do not by themselves give
a uniform bound on all solutions. In this example the degeneration
is generated by the conformal symmetry group. For general data,
excluding or controlling concentration requires additional arguments;
one cannot simply invoke elliptic compactness from (328.1).

\paragraph{What the obstruction does and does not characterize.}
The statement ``there exists a positive weight making the integrals
vanish'' is weaker than finding that weight as $e^{2u}$ for a
solution. The PDE links the weight to its own Laplacian, while a
finite list of moment equations does not. Conversely, evaluating
(328.3) with weight one and rejecting a $K$ on that basis is not
legitimate. The necessary conditions depend on the unknown metric.

Likewise local invertibility near $u=0$ on even functions cannot be
continued along an arbitrary path of curvatures without control of
the linearized operator and compactness. Along such a path solutions
can leave the small neighborhood, develop degeneracy, or concentrate.
Formula (328.7) gives an explicit reason that a bound must be proved
rather than inferred from smoothness of the right side.

\paragraph{Checks and outcome.}
The saved diagnostic differentiates the polynomial $P_2$ through
the zonal Laplacian $((1-z^2)P_2')'$, checks the coefficient in
(328.6), and numerically compares the radial mass formula (328.8)
with its antiderivative. These limited calculations do not certify
existence for an arbitrary curvature function.

The proved results are the exact conformal-vector obstruction,
the nonexistence of all nonconstant affine coordinate curvatures,
the small even-data existence theorem, and the explicit solution
families. The first general missing step is a complete criterion
that controls both obstruction and concentration for arbitrary
smooth data. The global characterization remains \textbf{UNRESOLVED}.
\subsection{Sources}
\begin{itemize}
\item[S1]Struwe, Bubbling and Topological Degeneration in the Calculus of Variations (2023).
\url{https://link.springer.com/article/10.1007/s00032-022-00370-1}.
\textit{Formulation source.}
\item[S2]Status source for Nirenberg Problem on the Two-Sphere.
\url{https://arxiv.org/abs/2603.29544}.
\textit{Further reference; not a proof endorsement.}
\end{itemize}
\end{document}
